\documentclass[11pt]{article}
% a typical CS/ML/physics preamble; maketitle prefetches the TeX Live files this needs (scripts: maketitle-lazy/base-trace)
\usepackage[utf8]{inputenc}
\usepackage[T1]{fontenc}
\usepackage{times}
\usepackage{geometry}
\usepackage{amsmath,amssymb,amsthm,amsfonts,mathtools,bm}
\usepackage{graphicx,xcolor,booktabs,multirow,nicefrac,microtype,url,enumitem,float,wrapfig}
\usepackage{caption,subcaption}
\usepackage[numbers]{natbib}
\usepackage{hyperref,cleveref}
\newtheorem{theorem}{Theorem}
\title{A Common Preamble}
\author{Anon}
\begin{document}
\maketitle
\begin{abstract}Short abstract with \emph{emphasis}, \textbf{bold} and $\alpha+\beta=\gamma$.\end{abstract}
\section{Intro}\label{sec:intro}
Text \cite{a} with math $\mathbb{R}^n$, $\mathcal{L}$, $\bm{x}$, $\nicefrac{1}{2}$ and \texttt{code}, \textsc{caps}, \textit{italic}. See \cref{sec:intro}.
\begin{equation}\label{eq:1} \sum_{i=1}^n x_i^2 \leq \int_0^1 f(x)\,dx \end{equation}
\begin{table}[h]\centering\begin{tabular}{lcc}\toprule A & B & C\\\midrule \multirow{2}{*}{x} & 1 & 2\\ & 3 & 4\\\bottomrule\end{tabular}\caption{T}\end{table}
\begin{theorem}A theorem.\end{theorem}\begin{proof}Trivial.\end{proof}
\begin{itemize}\item one\item two\end{itemize}
\url{https://example.org}
\bibliographystyle{plainnat}
\bibliography{refs}
\end{document}
